% !TeX root = ../Book1.tex
\begin{figure}[htbp]
	\centering
	\begin{tikzpicture}[scale=2.2,font=\small]
		\draw[thick,GeometryGray]
			(22.5:{1/cos(22.5)})--(67.5:{1/cos(22.5)})--(112.5:{1/cos(22.5)})--(157.5:{1/cos(22.5)})--(202.5:{1/cos(22.5)})--(247.5:{1/cos(22.5)})--(292.5:{1/cos(22.5)})--(337.5:{1/cos(22.5)})--cycle;
		\filldraw[fill=GeometryBlue!10,draw=GeometryBlue,thick]
			(0:1)--(45:1)--(90:1)--(135:1)--(180:1)--(225:1)--(270:1)--(315:1)--cycle;
		\draw[black,thick] (0,0) circle[radius=1];
		\foreach \a in {0,45,90,135,180,225,270,315}{
			\fill[GeometryBlue] (\a:1) circle[radius=0.018];
		}
		\draw[->,GeometryGray] (0,0)--(0:0.78) node[midway,below] {\(u_0\)};
		\node[GeometryBlue] at (-0.3,-0.1) {\(Q\)};
		\node[right] at (1.16,0.85) {\(P\)：支撑半平面的交};
		\draw[GeometryGray] (1.12,0.83)--(1,{tan(22.5)});
		\node[right] at (1.16,0.35) {\(\mathbb B\)：单位圆盘};
		\draw (1.12,0.33)--(0.95,0.31);
		\node[right,GeometryBlue] at (1.16,-0.18) {\(Q\)：支撑点的凸包};
	\end{tikzpicture}
	\caption[有限方向给出的内外近似]{八个等角单位方向给出 \(Q\subseteq\mathbb B\subseteq P\)。内多边形连接取到支撑值的点，外多边形保留相应支撑半平面；未采样方向上的两侧间隙由误差估计控制。}
	\label{g1:fig:support-polytope-approximation}
\end{figure}
